% =========================================================================
%  Guide de démonstration — Phase de démo soutenance RansomShield
%  Codjo Fréjus Loïc SANGRONIO — IFRI-UAC — Licence Sécurité Informatique
%  Durée totale : 4 minutes | 5 étapes | 3 VMs KVM
%  Compilateur : pdfLaTeX
% =========================================================================

\documentclass[a4paper,11pt]{article}
\usepackage[utf8]{inputenc}
\usepackage[T1]{fontenc}
\usepackage[french]{babel}
\usepackage{geometry}
\geometry{top=1.6cm,bottom=1.6cm,left=2cm,right=2cm}
\usepackage{xcolor}
\usepackage{tcolorbox}
\tcbuselibrary{skins,breakable}
\usepackage{fontawesome5}
\usepackage{booktabs}
\usepackage{tabularx}
\usepackage{array}
\usepackage{enumitem}
% titlesec remplacé par redéfinition directe
\usepackage{fancyhdr}
\usepackage[colorlinks=true,linkcolor=blue!60!black]{hyperref}
\usepackage{multicol}

% ── Couleurs ──────────────────────────────────────────────────────────────
\definecolor{ifriblue} {RGB}{  0,  82, 165}
\definecolor{darkblue} {RGB}{ 10,  40,  90}
\definecolor{accentred}{RGB}{220,  70,  30}
\definecolor{normgreen}{RGB}{ 22, 128,  65}
\definecolor{darkbg}   {RGB}{ 18,  30,  60}
\definecolor{lightblue}{RGB}{232, 240, 252}
\definecolor{warnbg}   {RGB}{255, 248, 220}
\definecolor{critbg}   {RGB}{254, 236, 236}

% ── Boîtes ────────────────────────────────────────────────────────────────

% Boîte de timing (bleue)
\newtcolorbox{etapebox}[2]{
  colback=ifriblue!8, colframe=ifriblue,
  fonttitle=\bfseries\normalsize\color{white},
  title={\faStopwatch\quad Étape #1 \hfill \textit{\normalfont\small #2}},
  boxrule=1.2pt, arc=5pt,
  left=7pt, right=7pt, top=5pt, bottom=5pt,
  breakable
}

% Boîte "À DIRE" (gris foncé italique)
\newtcolorbox{direbox}{
  colback=darkbg!92, colframe=darkbg,
  fontupper=\itshape\small\color{white!90!darkbg},
  boxrule=0pt, arc=3pt,
  left=8pt, right=8pt, top=4pt, bottom=4pt,
  title={\faCommentDots\quad \textbf{\normalfont\small À DIRE}},
  coltitle=white, fonttitle=\bfseries\small,
  attach boxed title to top left,
  boxed title style={colback=darkbg,arc=2pt,boxrule=0pt},
  breakable
}

% Boîte "À MONTRER / CLIQUER" (bleue claire)
\newtcolorbox{montrerbox}{
  colback=lightblue, colframe=ifriblue!70,
  boxrule=0.8pt, arc=3pt,
  left=6pt, right=6pt, top=3pt, bottom=3pt,
  title={\faMousePointer\quad \textbf{\normalfont\small À MONTRER / CLIQUER}},
  coltitle=white, fonttitle=\bfseries\small,
  attach boxed title to top left,
  boxed title style={colback=ifriblue,arc=2pt,boxrule=0pt},
  breakable
}

% Boîte terminal / commande (fond noir)
\newtcolorbox{cmdbox}{
  colback=darkbg, colframe=darkbg!60!black,
  fontupper=\ttfamily\small\color{normgreen!90!white},
  boxrule=0.8pt, arc=3pt,
  left=8pt, right=8pt, top=4pt, bottom=4pt,
  breakable
}

% Boîte alerte / plan de secours (orange)
\newtcolorbox{warnbox}[1]{
  colback=warnbg, colframe=accentred!70,
  fonttitle=\bfseries\small\color{accentred!80!black},
  title={\faExclamationTriangle\quad #1},
  boxrule=1pt, arc=3pt,
  left=6pt, right=6pt, top=4pt, bottom=4pt,
  breakable
}

% Boîte checklist (verte)
\newtcolorbox{checkbox}[1]{
  colback=normgreen!8, colframe=normgreen!70,
  fonttitle=\bfseries\small\color{normgreen!70!black},
  title={\faCheckSquare\quad #1},
  boxrule=1pt, arc=3pt,
  left=6pt, right=6pt, top=4pt, bottom=4pt,
  breakable
}

% ── En-têtes / pieds de page ──────────────────────────────────────────────
\pagestyle{fancy}
\fancyhf{}
\renewcommand{\headrulewidth}{0.5pt}
\fancyhead[L]{\small\color{gray} Guide de démonstration --- RansomShield}
\fancyhead[R]{\small\color{gray} Loïc SANGRONIO --- IFRI-UAC --- Sécurité Informatique}
\fancyfoot[C]{\small\thepage}

% ── Titres de sections ────────────────────────────────────────────────────
\makeatletter
\renewcommand{\section}{\@startsection{section}{1}{\z@}%
  {1.5ex \@plus 1ex \@minus .2ex}%
  {1ex \@plus .2ex}%
  {\large\bfseries\color{ifriblue}}}
\makeatother

\setlength{\parskip}{4pt}
\setlength{\parindent}{0pt}

% ── Macros utiles ─────────────────────────────────────────────────────────
\newcommand{\ok}{\textcolor{normgreen}{\faCheckCircle}\ }
\newcommand{\clic}[1]{\texttt{\small #1}}
\newcommand{\resultat}[1]{\textbf{\textcolor{ifriblue}{$\Rightarrow$ #1}}}

% =========================================================================
\begin{document}
% =========================================================================

% ── Page de titre ─────────────────────────────────────────────────────────
\begin{center}
  {\LARGE\bfseries\color{ifriblue} Guide de la Phase de Démonstration — 4 minutes}\\[6pt]
  {\Large\bfseries RansomShield}\\[4pt]
  {\large Système de détection et de réponse aux attaques ransomware}\\[8pt]
  \hrule height 1.5pt
  \vspace{4pt}
  {\small Loïc SANGRONIO \quad\textbar\quad IFRI-UAC
    \quad\textbar\quad Licence Sécurité Informatique
    \quad\textbar\quad 2023--2024}
  \vspace{4pt}
  \hrule height 1.5pt
\end{center}

\vspace{6pt}

% ── Vue d'ensemble ────────────────────────────────────────────────────────
\begin{tcolorbox}[colback=lightblue,colframe=ifriblue,arc=5pt,boxrule=1pt]
  \begin{center}
  \renewcommand{\arraystretch}{1.3}
  \begin{tabular}{@{}c l c l@{}}
    \toprule
    \textbf{\#} & \textbf{Ce qu'on montre} & \textbf{Durée} & \textbf{Résultat attendu} \\
    \midrule
    1 & Tableau de bord + agents actifs   & 40 s       & Vue globale de la console SOC \\
    2 & Attaque réelle depuis une VM      & 1 min      & Événements injectés en SSH \\
    3 & Détection en temps réel           & 50 s       & Alerte critique + incident auto \\
    4 & Réponse SOC + timeline d'audit    & 55 s       & Approbation + rollback + audit \\
    5 & Simulation kill chain complète    & 35 s       & 22 événements, pipeline validé \\
    \midrule
    & \textbf{Total} & \textbf{4 min} & \\
    \bottomrule
  \end{tabular}
  \end{center}
\end{tcolorbox}

% =========================================================================
\section{Préparation (5 minutes avant la démo)}
% =========================================================================

\begin{checkbox}{Ce qu'il faut vérifier avant de commencer}
\begin{itemize}[leftmargin=1.5em,itemsep=3pt]
  \item[\ok] Démarrer le serveur SOC~:
    \begin{cmdbox}
cd /home/sangronio/ransomshield/backend-laravel\\
php artisan serve --host=0.0.0.0 --port=8000
    \end{cmdbox}

  \item[\ok] Vérifier que les 3 VMs sont actives~:
    \begin{cmdbox}
virsh list --all
\# Résultat attendu : rs-client-1, rs-fileserver, rs-attacker en "running"
    \end{cmdbox}

  \item[\ok] Ouvrir le navigateur sur
    \clic{http://localhost:8000/console/dashboard} et se connecter~:
    \begin{itemize}[leftmargin=1em]
      \item E-mail~: \texttt{admin@ransomshield.local}
      \item Mot de passe + code 2FA TOTP (application Google Authenticator)
    \end{itemize}

  \item[\ok] Ouvrir deux fenêtres côte à côte~:
    fenêtre 1 = navigateur SOC \quad|\quad fenêtre 2 = terminal SSH.

  \item[\ok] Nettoyer les événements de test précédents si la console est bruyante
    (optionnel, pour une démonstration propre).
\end{itemize}
\end{checkbox}

\vspace{4pt}

\begin{warnbox}{Si un agent VM est hors ligne}
\begin{cmdbox}
ssh ubuntu@10.20.0.81     \# mot de passe : ransomshield2026\\
sudo systemctl restart ransomshield-agent\\
sudo systemctl status ransomshield-agent    \# doit afficher "active (running)"
\end{cmdbox}
Vérifier ensuite dans la console que le heartbeat se met à jour.
\end{warnbox}

\newpage

% =========================================================================
\section{Déroulé étape par étape}
% =========================================================================

% ─── ÉTAPE 1 ──────────────────────────────────────────────────────────────
\begin{etapebox}{1}{0:00 -- 0:40 \textnormal{(40 secondes)}}

\begin{montrerbox}
\begin{enumerate}[leftmargin=1.5em,itemsep=3pt]
  \item Aller sur \clic{/console/dashboard}.
  \item Montrer les compteurs en haut~: alertes actives, incidents en cours, agents critiques, actions en attente.
  \item Montrer le graphique d'événements (mis à jour toutes les 30\,s).
  \item Aller sur \clic{/console/agents}.
  \item Montrer les 3 VMs actives avec leur IP, heartbeat et niveau de risque~:\\
    \texttt{rs-client-1 (10.20.0.81)} \quad \texttt{rs-fileserver (10.20.0.38)} \quad \texttt{rs-attacker (10.20.0.72)}.
\end{enumerate}
\end{montrerbox}

\vspace{5pt}

\begin{direbox}
« Voici la console SOC de RansomShield — elle centralise la supervision de tous les terminaux.
Les compteurs en haut donnent l'état en temps réel : alertes actives, incidents, agents critiques, actions en attente.
Trois agents Python tournent sur trois machines virtuelles KVM en réseau isolé, chacun avec un heartbeat toutes les 30 secondes. »
\end{direbox}

\end{etapebox}

\vspace{6pt}

% ─── ÉTAPE 2 ──────────────────────────────────────────────────────────────
\begin{etapebox}{2 --- Attaque réelle depuis une VM}{0:40 -- 1:40 \textnormal{(1 minute)}}

\begin{montrerbox}
Basculer sur le terminal SSH (fenêtre 2) et jouer les commandes suivantes~:
\end{montrerbox}

\begin{cmdbox}
\# Se connecter à la VM victime\\
ssh ubuntu@10.20.0.81\\
\\
\# Créer un fichier document normal\\
touch /home/ubuntu/Documents/rapport\_demo.docx\\
sleep 2\\
\\
\# Chiffrement simulé : renommer avec extension .locked\\
mv /home/ubuntu/Documents/rapport\_demo.docx \textbackslash\\
\quad /home/ubuntu/Documents/rapport\_demo.docx.locked\\
sleep 2\\
\\
\# Déposer une note de rançon\\
echo "Your files are encrypted. Pay 1 BTC." \textbackslash\\
\quad > /home/ubuntu/HOW\_TO\_DECRYPT.txt
\end{cmdbox}

\vspace{4pt}

\begin{direbox}
« Je me connecte en SSH sur la VM victime et simule un chiffrement :
renommage avec l'extension \texttt{.locked}, puis dépôt d'une note de rançon.
L'agent Python surveille ce dossier via Watchdog et inotify.
Ces événements sont placés dans la file locale SQLite, puis transmis au SOC via l'API sécurisée. »
\end{direbox}

\end{etapebox}

\vspace{6pt}

% ─── ÉTAPE 3 ──────────────────────────────────────────────────────────────
\begin{etapebox}{3 --- Détection en temps réel}{1:40 -- 2:30 \textnormal{(50 secondes)}}

\begin{montrerbox}
Revenir sur le navigateur (fenêtre 1) --- ne pas rafraîchir manuellement.
\begin{enumerate}[leftmargin=1.5em,itemsep=3pt]
  \item \clic{/console/events} --- l'événement \texttt{file\_encrypted\_extension}
    de \texttt{rs-client-1} est apparu avec un score de \textbf{135} et le niveau \textbf{Critique}.
  \item \clic{/console/alerts} --- l'alerte critique a été générée automatiquement
    avec ses signaux (extension suspecte + note de rançon) et l'agent source.
  \item \clic{/console/incidents} --- l'incident a été ouvert automatiquement.
    Cliquer dessus pour afficher la fiche~: alertes rattachées, signaux, actions proposées.
\end{enumerate}
\end{montrerbox}

\vspace{5pt}

\begin{direbox}
« En moins de 5 secondes après le renommage sur la machine virtuelle,
le moteur de détection a calculé un score de 135 à partir des signaux cumulés~:
+80 pour l'extension suspecte, +55 pour la note de rançon.
Ce score dépasse le seuil critique configuré à 100.
Une alerte a été générée et un incident de sécurité a été ouvert automatiquement dans la console. »
\end{direbox}

\resultat{Score 135 → Niveau Critique → Alerte + Incident ouverts automatiquement en $< 5$\,s}

\end{etapebox}

\vspace{6pt}

% ─── ÉTAPE 4 ──────────────────────────────────────────────────────────────
\begin{etapebox}{4 --- Réponse SOC et timeline d'audit}{2:30 -- 3:25 \textnormal{(55 secondes)}}

\begin{montrerbox}
\begin{enumerate}[leftmargin=1.5em,itemsep=3pt]
  \item \clic{/console/protection-actions} --- montrer la file d'approbation~:
    actions proposées par le moteur (scan actif, isolation réseau).
  \item Cliquer \textbf{Approuver} sur l'action \textit{scan actif}~: l'agent reçoit la commande et exécute le scan.
  \item Si l'action ``isolation réseau'' est disponible~: l'approuver, puis montrer
    le bouton \textbf{Rollback} pour lever l'isolation.
  \item \clic{/console/incidents/\{id\}/timeline} --- afficher la timeline chronologique~:
    événement reçu $\to$ alerte créée $\to$ incident ouvert $\to$ action proposée $\to$ action approuvée.
\end{enumerate}
\end{montrerbox}

\vspace{5pt}

\begin{direbox}
« Dans RansomShield, le moteur propose — l'opérateur décide.
J'approuve l'action de scan : l'agent reçoit la commande et l'exécute.
Si j'avais isolé la machine, le bouton rollback permet de lever cette isolation en un clic.
Chaque décision est horodatée dans la timeline : de la détection à la réponse, tout est tracé. »
\end{direbox}

\resultat{Timeline complète~: de la détection à la décision, tout est tracé et horodaté}

\end{etapebox}

\vspace{6pt}

% ─── ÉTAPE 5 ──────────────────────────────────────────────────────────────
\begin{etapebox}{5 --- Simulation de la kill chain complète}{3:25 -- 4:00 \textnormal{(35 secondes)}}

\begin{montrerbox}
\begin{enumerate}[leftmargin=1.5em,itemsep=3pt]
  \item Aller sur \clic{/console/simulation}.
  \item Sélectionner l'agent cible~: \textbf{rs-client-1}.
  \item Choisir le scénario \textbf{``Kill chain complète''} (22 événements).
  \item Cliquer \textbf{Lancer la simulation}.
  \item Observer la cascade d'alertes sur \clic{/console/alerts}.
\end{enumerate}
\end{montrerbox}

\vspace{5pt}

\begin{direbox}
« Le simulateur injecte 22 événements dans le même pipeline que pour une vraie attaque.
Tous sont marqués \texttt{is\_simulation=true} pour un nettoyage facile.
Cela démontre la résilience du moteur face à un volume élevé d'événements, sans toucher l'environnement réel. »
\end{direbox}

\resultat{22 événements traités, aucun perdu~: la résilience de la file SQLite est démontrée}

\end{etapebox}

% =========================================================================
\section{Récapitulatif des capacités démontrées}
% =========================================================================

\begin{checkbox}{Ce que le jury a pu observer}
\renewcommand{\arraystretch}{1.2}
\begin{tabular}{@{}l l c@{}}
\toprule
\textbf{Capacité} & \textbf{Prouvée par} & \textbf{Étape} \\
\midrule
Multi-agents sur VMs réelles       & 3 agents actifs en console          & 1 \\
Surveillance en temps réel         & Événement détecté en $<$ 5\,s       & 2--3 \\
Analyse comportementale            & Score 135, niveau critique auto     & 3 \\
Incident auto-ouvert               & Fiche incident avec signaux         & 3 \\
Contrôle humain                    & File d'approbation (approuver/rejeter) & 4 \\
Rollback                           & Isolation réseau levée en un clic  & 4 \\
Timeline d'audit complète          & Chronologie horodatée              & 4 \\
Résilience du pipeline             & 22 événements sans perte            & 5 \\
\bottomrule
\end{tabular}
\end{checkbox}

% =========================================================================
\section{Plan de secours (si quelque chose échoue)}
% =========================================================================

\begin{warnbox}{La VM ne répond pas en SSH}
Injecter un événement directement via \texttt{curl} depuis le poste local~:
\begin{cmdbox}
curl -s -X POST http://127.0.0.1:8000/api/agent/events \textbackslash\\
  -H "Content-Type: application/json" \textbackslash\\
  -H "X-Agent-Secret: y0QUhoaW57aYsc0xd9tI7vbj2M18HogBDWPIAIe9xh2mTEo104Bv3vrNisdu" \textbackslash\\
  -d '\{"agent\_uuid":"96c75bc2-afcd-442f-8888-1f28b7b661ad",\\
        "event\_type":"file\_encrypted\_extension",\\
        "file\_path":"/home/ubuntu/rapport\_demo.docx.locked",\\
        "file\_extension":".locked","file\_name":"rapport\_demo.docx.locked"\}'
\end{cmdbox}
La réponse JSON confirme~: score, niveau de risque, alerte créée, incident ouvert.\\
Passer directement à l'\textbf{étape 5} (simulation) si les étapes 2--3 posent problème.
\end{warnbox}

\vspace{4pt}

\begin{warnbox}{Le navigateur ne montre pas les nouvelles alertes}
Forcer un rafraîchissement manuel (F5) ou aller directement sur
\clic{/console/events} et filtrer par niveau ``Critique''.
\end{warnbox}

\vspace{4pt}

\begin{warnbox}{La simulation ne se lance pas}
Vérifier qu'un agent \textbf{actif} (heartbeat $<$ 2 min) est sélectionné.
Utiliser l'agent local \texttt{DESKTOP-TEST01} en dernier recours.
\end{warnbox}

\vspace{6pt}
\begin{center}
\small\color{gray}
Document interne --- soutenance RansomShield\\
Loïc SANGRONIO --- IFRI-UAC --- 2023--2024
\end{center}

\end{document}
